
\subsection{非干净相交拉格朗日子流形的弗洛尔上同调}
\label{subsec:6-5-4}

现在利用定理~\ref{thm:6-1-25}，
为 $L^{(0)}$ 与 $L^{(1)}$
不一定干净相交的情形定义
$\Lambda_{0,\mathrm{nov}}$ 系数的弗洛尔上同调。

假设拉格朗日对子
$(L^{(0)},L^{(1)})$ 相对 spin，
并在 bulk 形变后弱无障碍。
取
$\mathbf b_j\in
\mathcal M_{\mathrm{weak,def}}(L^{(j)})$
（$j=0,1$），使
$\pi_{\mathrm{amb}}(\mathbf b_0)
=\pi_{\mathrm{amb}}(\mathbf b_1)$ 且
$\mathfrak{PO}(\mathbf b_0)
=\mathfrak{PO}(\mathbf b_1)$。
此时可以定义
\[
HF\bigl((L^{(1)},\mathbf b_1),
(L^{(0)},\mathbf b_0);
\Lambda_{\mathrm{nov}}\bigr),
\]
因为 $\Lambda_{\mathrm{nov}}$ 系数的
弗洛尔上同调在哈密顿同痕下不变；
总可扰动 $L^{(i)}$ 使其横截相交，
所得上同调典范同构。
我们希望借定理~\ref{thm:6-1-25}
定义 $\Lambda_{0,\mathrm{nov}}$ 系数的
$HF((L^{(1)},\mathbf b_1),
(L^{(0)},\mathbf b_0);
\Lambda_{0,\mathrm{nov}})$。

考虑两列哈密顿微分同胚
$\psi_m^{(0)},\psi_m^{(1)}$
（$m=1,2,3,\ldots$），满足
\begin{equation}
\psi_m^{(0)}(L^{(0)})
\text{ 与 }\psi_m^{(1)}(L^{(1)})
\text{ 干净相交。}
\tag{6.5.35.1}\label{eq:6-5-35-1}
\end{equation}
\begin{equation}
\lim_{m\to\infty}
\operatorname{dist}(\psi_m^{(i)},\mathrm{id})=0
\quad(i=0,1),
\tag{6.5.35.2}\label{eq:6-5-35-2}
\end{equation}
其中 $\operatorname{dist}$ 是霍费尔距离。
下述命题~\ref{prop:6-5-38}
声称，这时
\[
HF\bigl(
(\psi_m^{(1)}(L^{(1)}),
\psi_{m*}^{(1)}\mathbf b_1),
(\psi_m^{(0)}(L^{(0)}),
\psi_{m*}^{(0)}\mathbf b_0);
\Lambda_{0,\mathrm{nov}}\bigr)
\]
在下面定义的意义下随 $m\to\infty$ 收敛。

\begin{definition}{6.5.36}\label{def:6-5-36}
对 $m=1,2,\ldots,\infty$，
设 $\lambda_{m,\downarrow1},\ldots,
\lambda_{m,\downarrow b_m}$
是正数列，满足
$\lambda_{m,\downarrow i}
\ge\lambda_{m,\downarrow i+1}$。
假设有限 $m$ 时 $b_m<\infty$，
但允许 $b_\infty=\infty$；
在后种情形，另要求
$\lim_{i\to\infty}
\lambda_{\infty,\downarrow i}=0$。
若 $\lim_{m\to\infty}b_m=b_\infty$，
且对每个 $i$ 有
$\lim_{m\to\infty}
\lambda_{m,\downarrow i}
=\lambda_{\infty,\downarrow i}$，
则称
\[
\lim_{m\to\infty}
(\lambda_{m,\downarrow1},\ldots,
\lambda_{m,\downarrow b_m})
=
(\lambda_{\infty,\downarrow1},\ldots,
\lambda_{\infty,\downarrow b_\infty}).
\]
\end{definition}

\begin{remark}{6.5.37}\label{rem:6-5-37}
在命题~\ref{prop:6-5-38} 的情形，
可能有无穷多个 $\lambda_{\infty,\downarrow i}$，
且 $\lim_{i\to\infty}
\lambda_{\infty,\downarrow i}=0$。
此时需要按
$\lambda_{\infty,\downarrow i}
\ge\lambda_{\infty,\downarrow i+1}$
排列这些指数。
\end{remark}

\begin{proposition}{6.5.38}\label{prop:6-5-38}
设 $\psi_m^{(0)},\psi_m^{(1)}$
是满足式~\eqref{eq:6-5-35-1}、
\eqref{eq:6-5-35-2} 的
两列哈密顿微分同胚，并记
\[
\begin{aligned}
&HF\bigl(
(\psi_m^{(1)}(L^{(1)}),
\psi_{m*}^{(1)}\mathbf b_1),
(\psi_m^{(0)}(L^{(0)}),
\psi_{m*}^{(0)}\mathbf b_0);
\Lambda_{0,\mathrm{nov}}\bigr)\\
&\qquad\cong
(\Lambda_{0,\mathrm{nov}})^a
\oplus
\bigoplus_{i=1}^{b_m}
\frac{\Lambda_{0,\mathrm{nov}}}
{T^{\lambda_{m,\downarrow i}}
\Lambda_{0,\mathrm{nov}}}.
\end{aligned}
\]
则存在定义~\ref{def:6-5-36} 所述的
$\lambda_{\infty,\downarrow1},\ldots,
\lambda_{\infty,\downarrow b_\infty}$，
使
\[
\lim_{m\to\infty}
(\lambda_{m,\downarrow1},\ldots,
\lambda_{m,\downarrow b_m})
=
(\lambda_{\infty,\downarrow1},\ldots,
\lambda_{\infty,\downarrow b_\infty})
\]
按定义~\ref{def:6-5-36} 的意义成立。
\end{proposition}

原书称命题~\ref{prop:6-5-38}
是定理~\ref{thm:6-1-25}
的直接推论。
\translatornote{定理 6.1.25
只控制远离零的正挠指数，
未排除额外指数趋于零，
也未控制这类指数的数目；
故它本身不足以推出定义 6.5.36
要求的 $b_m$ 收敛及极限指数全为正。
源书未补充零附近挠项的论证，
命题及后续定义的一般情形保留这一缺口，
见疑误 A-0400。}

\begin{definition}{6.5.39}\label{def:6-5-39}
定义
\[
\begin{aligned}
&HF\bigl((L^{(1)},\mathbf b_1),
(L^{(0)},\mathbf b_0);
\Lambda_{0,\mathrm{nov}}\bigr)\\
&\quad:=
(\Lambda_{0,\mathrm{nov}})^a
\oplus
\bigoplus_{i=1}^{b_\infty}
\frac{\Lambda_{0,\mathrm{nov}}}
{T^{\lambda_{\infty,\downarrow i}}
\Lambda_{0,\mathrm{nov}}}.
\end{aligned}
\]
\end{definition}

注意，原书称
$\lim_{m\to\infty}
(\lambda_{m,\downarrow1},\ldots,
\lambda_{m,\downarrow b_m})$
无需取子列即可收敛；
因而定义~\ref{def:6-5-39}
的右端不依赖 $\psi_m^{(j)}$ 的选取。

第 5 章已在 $L^{(1)}$ 与 $L^{(0)}$
干净相交时证明定理~\ref{thm:G}。
原书用命题~\ref{prop:6-5-38}
与定义~\ref{def:6-5-39}
处理一般情形。
定义直接给出 (G.3)、(G.4)；
利用定理~\ref{thm:6-1-25} 的证明
可以检查 (G.5)。
由于后文不使用 (G.5)，
原书略去该项检查的细节。

\begin{example}{6.5.40}\label{ex:6-5-40}
设
\[
\begin{aligned}
X&=
\left\{
x+\sqrt{-1}\,e^{-1/|x|}
\sin\frac1{|x|}
\,\middle|\,0<|x|\le1
\right\}\cup\{0\}
\subset\mathbb C,\\
Y&=[-1,1]\times\{0\}\subset\mathbb C.
\end{aligned}
\]
\translatornote{源书实体页 403
在 $X$ 的集合式中写 $|x|\le1$，
但 $x=0$ 处的 $1/|x|$ 未定义；
紧接着却要求 $F(p)=0$ 且
$F(L^{(0)}\cap U)=X$。
这里显式加入由振幅衰减唯一确定的
连续延拓点 $0$，见疑误 A-0401。}
可找到拉格朗日子流形
$L^{(1)},L^{(0)}\subset T^2$、
点 $p\in L^{(1)}\cap L^{(0)}$、
$p$ 的邻域 $U$，
以及到其像的辛微分同胚
$F:U\to F(U)\subset\mathbb C$，
使 $F(p)=0$，并且
$F(L^{(0)}\cap U)=X$、
$F(L^{(1)}\cap U)=Y$。
见图~\ref{fig:6-5-2}。

容易看出
$HF((L^{(1)},\mathbf b_1),
(L^{(0)},\mathbf b_0);
\Lambda_{0,\mathrm{nov}})$
不是有限生成的；换言之，
$b_\infty=\infty$。
相反，当 $L^{(1)}$ 与 $L^{(0)}$
干净相交时，
该弗洛尔上同调是有限生成的。
高维中也可得到类似例子，
例如把这个例子与
$(T^{2n},T^n)$ 作直积。
\end{example}

\par\medskip\noindent\begin{minipage}{\textwidth}
\centering
\setlength{\unitlength}{0.72mm}
\begin{picture}(145,92)
\linethickness{0.42pt}
\multiput(7,55)(3,0){24}{\line(1,0){1.5}}
\multiput(7,91)(3,0){24}{\line(1,0){1.5}}
\multiput(7,55)(0,3){12}{\line(0,1){1.5}}
\multiput(79,55)(0,3){12}{\line(0,1){1.5}}
\put(10,72){\line(1,0){67}}
\qbezier(12,80)(15,65)(20,68)
\qbezier(20,68)(25,67)(29,75)
\qbezier(29,75)(32,80)(35,73)
\qbezier(35,73)(38,69)(40,72)
\qbezier(40,72)(42,75)(44,72)
\qbezier(44,72)(45,70)(46,72)
\qbezier(46,72)(47,74)(48,72)
\qbezier(48,72)(49,71)(50,72)
\qbezier(50,72)(51,73)(52,72)
\qbezier(52,72)(54,70)(57,75)
\qbezier(57,75)(60,79)(64,69)
\qbezier(64,69)(70,64)(75,80)
\put(70,81){$\scriptstyle L^{(0)}$}
\put(70,68){$\scriptstyle L^{(1)}$}
\linethickness{0.55pt}
\qbezier(46,26)(46,48)(93,48)
\qbezier(93,48)(140,48)(140,26)
\qbezier(140,26)(140,4)(93,4)
\qbezier(93,4)(46,4)(46,26)
\linethickness{0.35pt}
\qbezier(56,27)(56,44)(93,44)
\qbezier(93,44)(130,44)(130,27)
\qbezier(130,27)(130,14.2)(109,11.0)
\qbezier(87,10.1)(56,11.5)(56,27)
\qbezier(67,28)(67,40)(93,40)
\qbezier(93,40)(119,40)(119,28)
\qbezier(119,28)(119,20.6)(109,17.7)
\qbezier(87,16.2)(67,17.5)(67,28)
\linethickness{0.55pt}
\qbezier(84,28.8)(84,33)(93,33)
\qbezier(93,33)(102,33)(102,28.2)
\qbezier(105.2,27.4)(105.2,23)(93,23)
\qbezier(93,23)(80.8,23)(80.8,27.4)
\linethickness{0.35pt}
\put(105.5,27.2){$\scriptstyle L^{(0)}$}
\put(131,26){$\scriptstyle L^{(1)}$}
\multiput(87,8)(3,0){8}{\line(1,0){1.5}}
\multiput(87,16)(3,0){8}{\line(1,0){1.5}}
\multiput(87,8)(0,2){5}{\line(0,1){1}}
\multiput(109,8)(0,2){5}{\line(0,1){1}}
\put(65,54){\vector(3,-5){22}}
\end{picture}
\par\smallskip
{\renewcommand{\thefigure}{6.5.2}
\refstepcounter{figure}
{\normalfont\bfseries 图~\thefigure}\label{fig:6-5-2}}
\end{minipage}\par\medskip

利用定义~\ref{def:6-5-39}，
原书指出定理 J 与定理~\ref{thm:6-1-25}
也适用于 $L^{(1)}$ 与 $L^{(0)}$
不干净相交的情形。
结合例~\ref{ex:6-5-40}，
它给出下面的推论。

\begin{corollary}{6.5.41}\label{cor:6-5-41}
每个维数中都存在一对拉格朗日子流形
$L^{(1)},L^{(0)}$ 与函数 $N(\epsilon)$，
使 $\lim_{\epsilon\to0}N(\epsilon)=\infty$；
若哈密顿同痕 $\psi$ 到恒等映射的
霍费尔距离小于 $\epsilon$，则
\[
\#\bigl(\psi(L^{(0)})\cap L^{(1)}\bigr)
\ge N(\epsilon).
\]
\end{corollary}
